% !PRINTED 283
\weberpage{283}
%<<CONT
und erhalten eine Grundkurve

\weq{9}{\begin{aligned} x^4-y^4-z^4 &= 0,\\ x\,y\,z &= 2, \end{aligned}}

die nur vom 12. Grade ist. Wir haben dann weiter nach § 75:

\weq{10}{\begin{aligned} 4A &= x\xi-y\eta-z\varrho,\\ B &= x\xi y\eta+x\xi z\varrho-y\eta z\varrho, \end{aligned}}

und man erhält für den Grad $m$ der Korrespondenz nach
§ 75 (13):
\[
\begin{array}{ll}
n=3, & m=12,\\
n=11, & m=12,\\
n=19, & m=60,
\end{array}
\]
also wieder wie im vorigen Falle, wenn $n\equiv 0, 1 \pmod 3$ ist,
eine dreifach zu große Zahl.

Ist endlich $n\equiv 1 \pmod 4$, so setze man
\[
x=f^4(\omega),\ y=f_1^4(\omega),\ z=f_2^4(\omega)
\]
und erhält eine Grundkurve vom sechsten Grade.

Für $n=5$ ergibt sich die richtige Zahl $m=6$.

\noindent\rule{0.3\textwidth}{0.4pt}
